\scanpage{8}
\mspage{3}
\margin{The functor is fully faithful at best over a reduced base $S$ -- false for example over the dual numbers of a field (if $I\neq\unsure{\phi|\varnothing}$). To have a fully faithful functor, it is necessary in the \uline{proper} context, i.e. $X$ by $(\widehat X,I)$ (so that $X=\widehat X-I$)}
I call any element of the essential image a ``standard \inserted{relative} curve'' over $S$. The system $(\unsure{Y|X},\widetilde A,I,\sigma)$ which gives rise to it will be called the ``normalized compactification'' \inserted{(relative)} (by abuse of \deleted{\unsure{unscrewing|dividing} system} language, since it is not only the relative scheme $Y$, but $Y$ with the additional structure $\widetilde A,I,\sigma_{\widetilde A}\ldots$).

2. \uline{Graph associated with a standard curve}.

Let us return to the case of an algebraically closed base field $k$, to begin with. Let $\unsure{X|S}$ be a standard curve, hence $Y,\deleted{\unsure{\widetilde A,}}I,\widetilde A,\sigma_{\widetilde A}$.

Set
\[
(7)\qquad S=\pi_0(Y)\simeq\text{set of irreducible components of }X
\]
We then have the diagram of canonical maps between finite sets
\diagram{\((8)\qquad\)\begin{tikzcd}[ampersand replacement=\&]
\& I\arrow[d,"p"]\\
\widetilde A\arrow[r,"o"]\arrow[d,"q"'] \& S\\
A=\widetilde A/\sigma \&
\end{tikzcd}}
% Light lower-left pagination; hand unidentified.
\mspage{7}
